\subsection{Proof strategy}

The interpolation upper bound holds at every system size and every
finite depth. We reproduce it in \cref{upper:section}, including a
normalization that avoids assuming integrability of the multiplier
$\log a$. The lower bound has four steps.

\begin{enumerate}[label=\textup{\arabic*.},leftmargin=*]
\item \emph{Select a cavity reservoir.}
A reservoir is a common Gibbs system on $N$ spins used to compare
systems of sizes $N$ and $N+1$.
For bounded interaction and field laws, add bounded Poisson perturbations whose
number is sublinear in the system size. Their effect on the pressure
vanishes. At each fixed hierarchy depth, select system sizes and
deterministic perturbation parameters for which the difference between
successive expected logarithms after the auxiliary power-mean recursion
is at most an arbitrarily small error above the original pressure liminf. The same choices give concentration of every
perturbation score (\cref{cavity:section,pert:selection}).

\item \emph{Prescribe branching patterns.}
An elementary calculus for finite probability trees expands insertions
into signed sums over added replica paths. Independent marks select any
prescribed branching pattern. Score concentration then gives identities
for these signed sums (\cref{tree:section,pert:tree-identity}).

\item \emph{Concentrate all multioverlaps.}
Move each internal branching depth of a target tree one step later.
The first fresh split at each new depth supplies the same small factor
as the normalizing coefficient. A relative bound on the sum of absolute
coefficients preserves control of errors after division. Conditional
covariance then decays by induction on the number of replicas.
A second identity concentrates the remaining conditional mean
(\cref{decor:section}).

\item \emph{Recover independent cavity messages.}
Encode a uniformly chosen reservoir site's entire finite hierarchy by
successive conditional laws. Multioverlap concentration permits different
site labels to be replaced by independent copies of this process in
the cavity insertion. This yields an admissible trial value close to
the selected increment (\cref{comp:section}).
\end{enumerate}

The estimates in Step~3 are averages over separated interior branching
depths. The insertion expansion assigns vanishing probability to the
exceptional trees, so these averages are sufficient for Step~4. No
uniform concentration in the depth is needed. The resulting lower bound
holds for every bounded interaction law invariant under permutations of
its inputs, without factorization or positivity. Finally, normalization
and estimates uniform over all trial laws remove boundedness and give
\cref{main:theorem} under the stated first moment assumption
(\cref{reduction:section}).
